\subsection{Расчёт трудоёмкости и определение исполнителей}

В данном разделе производится оценка трудоёмкости работ по разработке программного средства обнаружения аномалий в Kubernetes-кластере на основе анализа audit-логов. Общие трудозатраты по внедрению и разработке $Q_p$ вычисляются по формуле:

\begin{equation}
Q_p = \sum_i T_i,
\end{equation}

где $T_i$ -- трудоёмкость $i$-й работы.

Этапы, описанные ранее, были разбиты на отдельные работы. Трудоёмкость работ была определена методом экспертных оценок. Трудоёмкость $T_i$ работы вычисляется по формуле:

\begin{equation}
T_i = \frac{3T_{i\min} + 2T_{i\max}}{5},
\end{equation}

где $T_{i\min}$ и $T_{i\max}$ -- минимальная и максимальная экспертная оценка трудоёмкости работы $i$, 

Каждый этап технологического процесса содержит определённый набор работ. Для оценки трудоёмкости каждой работы обычно используется метод экспертных оценок, в результате применения которого получают значения минимального $T_{i\min}$ и максимального $T_{i\max}$ времени выполнения каждой работы. События и работы представлены в таблице \ref{tab:tasks_tруд}

{
\fontsize{10}{10}\selectfont
\setlength{\tabcolsep}{3pt}

\begin{longtable}{|c|c|p{4cm}|p{1.6cm}|p{1.6cm}|p{1.5cm}|p{1.6cm}|p{2.5cm}|}
    \caption{Перечень работ} \label{tab:tasks_tруд} \\
    \hline
    Этап & № работы & Содержание работы & $T_{i\min}$ (чел/час) & $T_{i\max}$ (чел/час) & $T_i$\newline(чел/час) & $T_i$\newline(чел/дни) & Трудовые ресурсы \\
    \hline
    \endfirsthead

    \caption*{Продолжение таблицы \ref{tab:tasks_tруд}}\\ \hline
    Этап & № работы & Содержание работы & $T_{i\min}$ (чел/час) & $T_{i\max}$ (чел/час) & $T_i$\newline(чел/час) & $T_i$\newline(чел/дни) & Трудовые ресурсы \\ \hline
    \endhead

    \hline
    \endfoot

    \hline
    \endlastfoot

    \multirow{4}{*}{1}
        & 1  & Анализ предметной области и существующих решений
        & 40 & 80  & 56.0  & 7.0  & Аналитик \\ \cline{2-8}
        & 2  & Анализ audit-логов Kubernetes-кластера
        & 32 & 64  & 44.8  & 5.6  & Аналитик \\ \cline{2-8}
        & 3  & Формирование требований к системе
        & 24 & 56  & 36.8  & 4.6  & Аналитик \\ \cline{2-8}
        & 4  & Постановка задачи
        & 16 & 40  & 25.6  & 3.2  & Аналитик \\ \hline

    \multirow{3}{*}{2}
        & 5  & Проектирование общей архитектуры системы
        & 40 & 80  & 56.0  & 7.0  & Инженер \\ \cline{2-8}
        & 6  & Проектирование структуры данных и хранилища логов
        & 32 & 72  & 48.0  & 6.0  & Инженер \\ \cline{2-8}
        & 7  & Проектирование взаимодействия модулей
        & 24 & 56  & 36.8  & 4.6  & Инженер \\ \hline

    \multirow{3}{*}{3}
        & 8  & Подготовка и предобработка данных
        & 48 & 96  & 67.2  & 8.4  & Инженер \\ \cline{2-8}
        & 9  & Разработка моделей НС для обнаружения аномалий
        & 80 & 160 & 112.0 & 14.0 & Инженер \\ \cline{2-8}
        & 10 & Обучение моделей НС
        & 64 & 128 & 89.6  & 11.2 & Инженер \\ \hline

    \multirow{2}{*}{4}
        & 11 & Оценка качества моделей
        & 40 & 80  & 56.0  & 7.0  & Аналитик \\ \cline{2-8}
        & 12 & Сравнение моделей и выбор оптимальной
        & 24 & 56  & 36.8  & 4.6  & Аналитик \\ \hline

    \multirow{3}{*}{5}
        & 13 & Реализация модуля сбора и обработки логов
        & 64 & 120 & 86.4  & 10.8 & Инженер \\ \cline{2-8}
        & 14 & Реализация модуля анализа и обнаружения аномалий
        & 80 & 160 & 112.0 & 14.0 & Инженер \\ \cline{2-8}
        & 15 & Реализация вспомогательных сервисов и API
        & 48 & 96  & 67.2  & 8.4  & Инженер \\ \hline

    \multirow{2}{*}{6}
        & 16 & Интеграция модели НС в систему
        & 40 & 80  & 56.0  & 7.0  & Инженер \\ \cline{2-8}
        & 17 & Настройка взаимодействия компонентов системы
        & 32 & 64  & 44.8  & 5.6  & Инженер \\ \hline

    \multirow{3}{*}{7}
        & 18 & Функциональное тестирование
        & 64 & 120 & 86.4  & 10.8 & Инженер \\ \cline{2-8}
        & 19 & Нагрузочное тестирование
        & 40 & 80  & 56.0  & 7.0  & Инженер \\ \cline{2-8}
        & 20 & Исправление ошибок и оптимизация
        & 64 & 128 & 89.6  & 11.2 & Инженер \\ \hline

    \multirow{2}{*}{8}
        & 21 & Подготовка пользовательской документации
        & 32 & 72  & 48.0  & 6.0  & Аналитик \\ \cline{2-8}
        & 22 & Подготовка эксплуатационной документации
        & 32 & 72  & 48.0  & 6.0  & Аналитик \\ \hline

    \multicolumn{5}{|r|}{Итого $Q_p$:} 
        & 1360 
        & 170
        & \\ \hline

\end{longtable}
}

Таким образом, срок выполнения всех работ составил 1360 чел./часов или 170 чел./дней. 

Среднее число исполнителей, необходимое для выполнения всех этапов разработки в сроки,можно определить из соотношения:

\begin{equation}
N = \left\lceil \frac{Q_p}{F} \right\rceil,
\end{equation}

$F$ -- фонд рабочего времени.

Фонд рабочего времени вычисляется по формуле:

\begin{equation}
F = t \cdot \sum_i D_i,
\end{equation}

где $t$ -- количество часов в рабочем дне, 

$D_i$ -- количество рабочих дней в рассматриваемом периоде.

Так как разработка проекта планируется в течение 6 месяцев, количество рабочих дней в этом промежутке составляет 130 дней. При восьмичасовом рабочем дне общий фонд рабочего времени составит:

\begin{equation}
F = 8 \cdot 130 = 1040 \ \text{ч}.
\end{equation}

Тогда численность исполнителей равна:

\begin{equation}
N = \left\lceil \frac{1360}{1040} \right\rceil = 2.
\end{equation}

Следовательно, для разработки программного средства обнаружения аномалий в Kubernetes-кластере на основе анализа audit-логов необходимо участие двух исполнителей. Привлечённые специалисты также будут задействованы в последующем внедрении и сопровождении проекта.

Исполнителями проекта должны быть универсальные инженеры-программисты с опытом разработки программных решений в области анализа данных и информационной безопасности. От исполнителей требуется знание принципов работы Kubernetes, понимание архитектуры контейнеризированных приложений, владение методами анализа журналов событий, знание основ статистики и машинного обучения, умение разрабатывать программную документацию, а также владение английским языком на уровне, достаточном для работы с технической документацией.